% Generated draft: requires mathematical and source review.
\documentclass[11pt]{article}
\usepackage[margin=25mm]{geometry}
\usepackage{fontspec}
\setmainfont{Latin Modern Roman}
\usepackage{amsmath,amssymb,mathtools,unicode-math}
\setmathfont{Latin Modern Math}
\usepackage{enumitem,needspace,xurl,hyperref}
\setlength{\parindent}{0pt}
\setlength{\parskip}{5pt}
\setlength{\emergencystretch}{3em}
\newcommand{\N}{\mathbb{N}}
\newcommand{\Z}{\mathbb{Z}}
\newcommand{\Q}{\mathbb{Q}}
\newcommand{\R}{\mathbb{R}}
\newcommand{\C}{\mathbb{C}}
\newcommand{\F}{\mathbb{F}}
\newcommand{\A}{\mathbb{A}}
\newcommand{\PP}{\mathbb{P}}
\newcommand{\E}{\mathbb{E}}
\newcommand{\eps}{\varepsilon}
\newcommand{\dd}{\,\mathrm{d}}
\begin{document}
\section*{UnsolvedMath ID: 30000473}
\textbf{Source problem code:} catalogue entry
\subsection{Definitions and mathematical statement}
\textbf{Conjecture 6, as reconstructed in the supplied record.} Suppose \(\lambda\) is a limit of supercompact cardinals, \(X\subseteq V_{\lambda+1}\), and \(j:L(X,V_{\lambda+1})\to L(X,V_{\lambda+1})\) is a proper elementary embedding. If square roots \(k_1,k_2\) of \(j\) agree on \(V_\lambda\) and both map \(L_\omega(X,V_{\lambda+1})\) to itself, must they agree below \(\Theta^{L(X,V_{\lambda+1})}\)? All these hypotheses are retained; in particular, no existence of such an embedding or roots is asserted.\par Write \(M=L(X,V_{\lambda+1})\) for the intended relative constructible universe, and \(V_\alpha\) for the cumulative hierarchy. A limit of supercompact cardinals conventionally means that supercompact cardinals below \(\lambda\) are unbounded in \(\lambda\), in the intended ambient universe. Elementarity means preservation of every first-order formula of the specified model language. The agreement hypothesis is \(k_1\upharpoonright V_\lambda=k_2\upharpoonright V_\lambda\). The source does not define the technical adjective \emph{proper} in this embedding context; it cannot be replaced merely by nontriviality or by being a proper-class map.\par The square-root operation is also unspecified. Composition roots satisfying \(k_i\circ k_i=j\) and roots under an application operation on elementary embeddings are different notions in rank-into-rank settings. No composition equation is silently assigned to the source term. Likewise, \emph{map \(L_\omega\) to itself} could mean preservation of its image or a specified fixed-object condition; it does not without further definition mean that every element is fixed. The exact construction of the stage \(L_\omega(X,V_{\lambda+1})\), the model language, and whether the roots themselves satisfy a properness condition must be recovered.\par The ordinal \(\Theta^M\) is not defined in the excerpt. In a higher-rank analogue one often uses the supremum of ordinals that are surjective images of \(V_{\lambda+1}\) by maps belonging to \(M\), whereas a determinacy-model \(\Theta\) based on real surjections is a different convention. The source must specify which is intended. Under an ordinal-action interpretation, \emph{agree below \(\Theta^M\)} means \(k_1(\alpha)=k_2(\alpha)\) for every ordinal \(\alpha<\Theta^M\); agreement on all sets of rank below \(\Theta^M\) is stronger and is not silently substituted. These definitions are essential to a fully determined conjecture.\par The supplied triage checked 2026-08-21 labels the item open while calling its extraction an ADR-like definitional and programmatic fragment. The later verification field reports direct recovery of all hypotheses and conclusion of Conjecture 6. That source-layer discrepancy and the missing technical definitions remain explicit. No current resolution or independent source verification is asserted.
\subsection{Short English statement}
The source asks for weak uniqueness of two embedding roots from their agreement on \(V_\lambda\) and a specified low-stage preservation condition. Properness, the root operation, \(\Theta\), and the meaning of agreement below it are technical missing definitions, so they are not guessed from ordinary composition or rank conventions.
\subsection{Sources}
\begin{itemize}\item Local UnsolvedMath registry, permanent ID 30000473, code catalogue entry; \emph{Weak Uniqueness of Elementary-Embedding Square Roots at \(\lambda\)}, Oberwolfach Reports Open Problems; source verification identifies Conjecture 6.\item \emph{Feinstrukturtheorie und Innere Modelle (2007)}, Oberwolfach Reports, DOI 10.4171/owr/2006/20, pp. 1215--1244. \url{https://doi.org/10.4171/owr/2006/20}. Proposed year and year tag: 2006; no proposer named.\item Supplied literature label \emph{original OWR report 20/2006}: \url{https://doi.org/10.4171/owr/2006/20}. No separate authoritative definitions of properness, roots, or \(\Theta\) are supplied.\item Literature triage checked 2026-08-21, status \texttt{open} with a fragment qualification; record updated 2026-09-16. The \texttt{corrected\_verified} field reports primary-report inspection and recovery of Conjecture 6, not independently checked here.\end{itemize}
\end{document}
